\clearpage
\fancyhead[L]{\small COMET.}
\SourceLeaf{425}
\begin{center}
\Large COMET.

\medskip
Göttingische gelehrte Anzeigen. No. 83. pp. 825..829. 24 May 1819.
\end{center}

\par\indent In issue 28 of this journal [p. 417 of this volume], we reported on the two comets of last year and already pointed out the likelihood that one of this comet was the same as the first of 1805. Professor \textsc{Encke}’s later calculations have increased this likelihood so greatly that there can hardly be any doubt about their identity. We have no hesitation in counting this addition to our knowledge of the solar system among the most remarkable discoveries of the nineteenth century, and we give the main points here from some letters that astronomer wrote to Councillor \textsc{Gauss}. Since no further usable observations of the comet of 1818 have yet become known beyond those printed in the issue cited, Professor \textsc{Encke} thought the best course for the time being was to examine the observations of the comet of 1805 afresh. The surprising result is that an ellipse with an orbital period of 3.42 years fits these observations best. For the comet of 1818, he had found an orbital period of 3.59 years from the elements cited. Both figures naturally still carried some uncertainty. Since 13.19 years have passed between the two observed times of passage through perihelion, the most likely conclusion, assuming the comets were identical, was that the comet had completed four revolutions in the interval (No one familiar with the subject can be surprised that the comet, though not among the very faintest, could only be seen through a telescope and went unnoticed on its three returns to perihelion in 1809, 1812, and 1815). With the orbital period adjusted to accord with this conclusion, Professor \textsc{Encke} calculated new elliptical elements from the observations of 1805. The new elements represented those observations most satisfactorily. It was a pleasant surprise that the individual elements of this new orbit came very much closer to those of the orbit calculated from the observations of 1818. Encouraged by this result, Mr. \textsc{Encke} decided to use Councillor \textsc{Gauss}’s method to calculate the perturbations those elements must have undergone over the whole 13 years under the influence of Jupiter. This laborious work too had its reward, for the agreement improved still further. We give the elements as calculated by Mr. \textsc{Encke} for the beginning of 1819 after allowing for Jupiter’s effects. They are based solely on the observations of 1805, with only the major axis set so that after four revolutions the time of passage through perihelion coincided with the one observed in 1819:

\begin{center}
\begin{tabular}{l l}
\multicolumn{2}{l}{Passage through perihelion, 1819 Jan. 27.27} \\
Longitude of perihelion . . . . . . . . . . . . & $156^\circ$ $59'$ $30''$ \\
Longitude of the ascending node . . . . . . . . . & 334\ \ 31\ \ \ 0 \\
\multicolumn{2}{l}{\quad both reckoned from the mean equinox of 1819} \\
Inclination of the orbit . . . . . . . . . . . . . . & 13\ \ 36\ \ 35 \\
Eccentricity angle . . . . . . . . . . . . . & 58\ \ \ 2\ \ 58 \\
Logarithm of the semimajor axis . . . . . . . . & 0.3451979 \\
\end{tabular}
\end{center}

\Needspace{12\baselineskip}\par\indent Finally, Mr. \textsc{Encke} determined the orbit anew from the observations \SourceLeaf{426}of 1818..1819 alone, assuming the logarithm of the semimajor axis $= 0.345$. He found that the observations could be represented much better this way than by the orbit calculated earlier. The new elements are as follows:

\begin{center}
\begin{tabular}{l l}
Time of perihelion, 1819 January . . . . . . . . & 27.27545 Seeberg time \\
Longitude of perihelion . . . . . . . . . . . . & $157^\circ\,5'\,53$: \\ 
Longitude of the node . . . . . . . . . . . . . . & $334\;\;43\;\;37$ \\
Inclination of the orbit . . . . . . . . . . . . . . & $13\;\;38\;\;42$ \\
Eccentricity angle . . . . . . . . . . . . . & $58\;\;6\;\;45.52$ \\
\end{tabular}
\end{center}

\par\indent These elements agree so closely with those derived from the observations of 1805 that the probability that the two comets are identical is equal to what is called certainty in matters outside pure mathematics.

\par\indent We might still wish to know precisely how far this splendid agreement depends on the assumed number of \emph{four} revolutions. It seems important to us that the orbits also be determined on the assumption of only three revolutions. One can foresee that the observations of 1805 and those of 1818 can both still be represented quite well by ellipses with an orbital period of 4.4 years. According to a preliminary estimate by Councillor \textsc{Gauss}, however, the agreement between the two orbits will be considerably poorer than between those given above. We hope Mr. \textsc{Encke}, who has already earned such great credit by his calculations so far, will undertake this investigation as well, and we are inclined to believe. that the result will allow us to judge the inadmissibility of assuming any number of revolutions other than four, even without calculating the perturbations in full. This investigation would become unnecessary if the comet could be observed again this year. In our northern regions there is admittedly little or no hope of that. In Italy, however, it does not seem quite impossible if instruments with very high light-gathering power are used. To help with such attempts, we therefore give an ephemeris calculated by Mr. \textsc{Encke} for the coming months. The times are $4^{\mathrm{u}}\,26^{\mathrm{m}}$ minutes Seeberg time.
\nopagebreak[4]
\begin{center}
\begin{tabular}{c|c|c|c}
1819 & Right & Southern & Log. of dist. \\
 & ascension & declination & from Earth \\
\hline
May\quad 31 & $336^\circ\,28'$ & $19^\circ\,43'$ & 0.2277 \\
June\quad 10 & $335\;\;27$ & $20\;\;22$ & 0.2168 \\
\phantom{June\quad} 20 & $333\;\;30$ & $21\;\;19$ & 0.2070 \\
\phantom{June\quad} 30 & $330\;\;36$ & $22\;\;28$ & 0.2004 \\
July\quad 10 & $326\;\;49$ & $23\;\;43$ & 0.1991 \\
\phantom{July\quad} 20 & $322\;\;26$ & $24\;\;54$ & 0.2049 \\
\phantom{July\quad} 30 & $317\;\;47$ & $25\;\;51$ & 0.2189 \\
Aug.\quad 9 & $313\;\;17$ & $26\;\;30$ & 0.2408 \\
\phantom{Aug.\quad} 19 & $309\;\;20$ & $26\;\;47$ & 0.2691 \\
\phantom{Aug.\quad} 29 & $306\;\;11$ & $26\;\;48$ & 0.3020 \\
Sept.\quad 8 & $303\;\;54$ & $26\;\;35$ & 0.3374 \\
\phantom{Sept.\quad} 18 & $302\;\;29$ & $26\;\;13$ & 0.3737 \\
\end{tabular}
\end{center}

\Needspace{7\baselineskip}\par\indent It should also be noted that this ephemeris was calculated from the earlier elements. The following two positions calculated from the newer elements give an idea of how far the newer elements differ:

\begin{center}
\begin{tabular}{c|c|c|c}
July\quad 30 & $320^\circ\,32$ & $24^\circ\,36'$ & 0.2097 \\
Aug.\quad\;\;9 & $315\;\;59$ & $25\;\;21$ & 0.2294 \\
\end{tabular}
\end{center}

\SourceLeaf{427}\par\indent We close this report in the hope that this most remarkable discovery will spur friends of astronomy to redouble their zeal in searching for comets, for everything suggests that it is only the beginning of an immense harvest that will ripen little by little.
